% FuzzPilot — Paper 1 arXiv preprint
%
% Build: pdflatex paper01.tex && bibtex paper01 && pdflatex paper01.tex && pdflatex paper01.tex
% Target: arXiv submission, cs.SE primary, cs.CR cross-list, CC-BY 4.0.
% Template choice: vanilla article (per user decision in plan
% /Users/qiaozhiyi/.claude/plans/lovely-wobbling-lovelace.md). Port to a
% conference template at Paper 3 time.

\documentclass[11pt]{article}

\usepackage[margin=1in]{geometry}
\usepackage{microtype}
\usepackage{booktabs}
\usepackage{graphicx}
\usepackage{xcolor}
\usepackage{hyperref}
\usepackage{tikz}
\usetikzlibrary{arrows.meta,positioning,shapes.geometric}
\usepackage{listings}
\usepackage{url}
\usepackage{cite}

\hypersetup{
  colorlinks=true,
  linkcolor=black,
  urlcolor=blue,
  citecolor=black,
}

\title{FuzzPilot: Off-Hot-Path LLM Control for AFL++ via\\
  Recipe-Guided Mutation and Micro-Campaign Validation}

\author{Anonymous Authors\thanks{To be de-anonymized at submission.}}

\date{\today}

\begin{document}
\maketitle

\begin{abstract}
% TODO Day 7: replace X\% with measured microbench number from
% results/paper01_ai_recipe_mutation/aggregated/microbench.csv (fp-active
% vs vanilla). Per outline: off-hot-path MUST appear in sentence 1.
Recent work places large language models inside the mutation loop of
greybox fuzzers, paying a steep throughput cost for every model call.
We argue that the model belongs \emph{off the hot path}: invoked only
at plateau boundaries to propose strategy, never to mutate bytes.
We present FuzzPilot, an AFL++ controller built around three ideas:
(i) a \textbf{recipe abstraction} that compiles model proposals into
structured guidance consumed by a native AFL++ custom mutator;
(ii) a \textbf{micro-campaign protocol} that validates each proposal
in a short isolated AFL++ run before promoting it; and
(iii) an \textbf{audit trail} of blackboards, agent decisions, and
recipe outcomes that makes every intervention reproducible after the
fact. On cJSON, FuzzPilot recovers from coverage plateaus faster than
AFL++ alone while keeping the mutator's native exec/sec within X\% of
vanilla AFL++. We release the system, configurations, and decision
logs to support reproduction.
\end{abstract}

% =====================================================================
\section{Introduction}
\label{sec:intro}
% TODO Day 7. Source: docs/papers/paper01_arxiv_placeholder.md
%   "Contributions" section — list the five contributions verbatim,
%   then say RQ on multi-target generalization, full ablation matrix,
%   and crash dedup are deferred to Paper 3 (full evaluation).
%   Off-hot-path framing MUST appear before contributions list.
% Page budget: 1.0

% =====================================================================
\section{Background and Threat to Throughput}
\label{sec:bg}
% TODO Day 7. Source: paper01_arxiv_placeholder.md, "Background &
%   Threat to Throughput" entry. Reuse the framing: AFL++ does
%   100k--1M exec/s; every per-input LLM call is 100--1000x slower.
%   Therefore any LLM-in-hot-path design caps throughput at model rate.
% Page budget: 0.5

% =====================================================================
\section{FuzzPilot Design}
\label{sec:design}
% TODO Day 2. Source: outline §"Sentence skeletons for the three pillars"
%   (already written verbatim in paper01_arxiv_placeholder.md, just port).
%   Anchor on Figure 1 (architecture). Implementation refs:
%   src/controller/run.cpp:411--461 (ablation switches reveal the design
%   joins).
% Page budget: 2.0

\begin{figure}[t]
  \centering
  % TODO Day 2 + Day 8. TikZ source in figures/F1_architecture.tikz.
  % Boxes: AFL++ runner, telemetry, plateau detector, blackboard, agents,
  % recipe store, native mutator. Edges show direction of data flow and
  % which edges cross the "hot path" boundary (the model never crosses).
  \fbox{\parbox{0.8\textwidth}{\centering\textit{[F1: FuzzPilot architecture --- placeholder]}\\TikZ in \texttt{figures/F1\_architecture.tikz}}}
  \caption{FuzzPilot architecture. The model sits off the hot path,
  invoked only by the plateau detector. The recipe store is the only
  channel into the native AFL++ custom mutator.}
  \label{fig:arch}
\end{figure}

% =====================================================================
\section{Recipe Abstraction and Native Mutator}
\label{sec:recipe}
% TODO Day 3. Source: mutators/fuzzpilot/ headers for the op table.
%   Cover: recipe schema, op vocabulary (insert/replace/splice/dict/
%   byte-range), dispatch cost (paid once per recipe load, not per
%   call), and the contract that the mutator never re-enters the model.
% Page budget: 1.0

\begin{figure}[t]
  \centering
  % TODO Day 3. TikZ source in figures/F2_recipe_lifecycle.tikz.
  % proposal -> schema validate -> micro-campaign -> reward -> promote/discard.
  \fbox{\parbox{0.8\textwidth}{\centering\textit{[F2: Recipe lifecycle --- placeholder]}\\TikZ in \texttt{figures/F2\_recipe\_lifecycle.tikz}}}
  \caption{Recipe lifecycle: proposal $\rightarrow$ schema validation
  $\rightarrow$ micro-campaign $\rightarrow$ reward $\rightarrow$ promote
  or discard.}
  \label{fig:lifecycle}
\end{figure}

% =====================================================================
\section{Plateau Detection and Micro-Campaign Protocol}
\label{sec:plateau}
% TODO Day 4. Source: src/controller/run.cpp:411--461 (ablation
%   switches: no-plateau, no-microcampaign, no-static-analysis,
%   no-mutator, random-recipe, random-reward, edges-only).
%   docs/papers/paper01_ghidra_and_x86_audit.md §1.4 explains how
%   Ghidra feeds the prompt (NOT the mutator) --- relevant for the
%   E2c (no-static-analysis) ablation interpretation.
% Page budget: 1.0

% =====================================================================
\section{Preliminary Evaluation}
\label{sec:eval}
% TODO Day 5. After aggregate + plots land. Source files:
%   results/paper01_ai_recipe_mutation/aggregated/{t1_per_run_summary.csv,
%   microbench.csv, plateau_recovery.csv, coverage_timeseries.csv}
%   results/paper01_ai_recipe_mutation/figures/{F3,F4,F5}.pdf
% Structure: RQ1 (plateau recovery via F3+F4), RQ2 (throughput via F5),
%   RQ3 (promotion rate + decision trail via T1).
% Acceptance gate footnote: cite paper01_preprint.yaml acceptance block.
% Page budget: 1.5

\subsection{Setup}
% TODO Day 5: target, modes, repeats, budget, hardware, model
%   (deepseek-chat at temperature 0.0), seed corpus, AFL++ commit,
%   FuzzPilot tag, Docker image digest.

\subsection{RQ1 --- Plateau recovery on cJSON}
% Figures: F3 (coverage curves), F4 (time-to-recover boxplot)
% Tables: T1 (per-run summary)

\subsection{RQ2 --- Mutator throughput}
% Two pieces of evidence:
%   (a) End-to-end (canonical): E1a vs E1b fuzzer_stats:execs_per_sec
%       median ratio, from aggregated/throughput_parity.csv. Manifest
%       gate: ratio >= 0.85. THIS is the real RQ2 number.
%   (b) Microbench (F5, supplementary): dispatch-symmetric microbench
%       where vanilla = libvanilla_havoc.so (AFL++ havoc-equivalent
%       through the same dlopen path as FuzzPilot mutator). Two gates:
%       fp-active vs fp-empty <= 5% (recipe is free, this is a real
%       paper claim) and fp-empty vs vanilla <= 5x (sanity; FuzzPilot
%       mutator does telemetry-ring writes + recipe-cache lookups
%       even with empty store, so a ~3x dispatch cost over a minimal
%       AFL-style havoc baseline is expected).
%   Write (a) first, then F5 as supplementary panel.
%   File refs: aggregated/throughput_parity.csv (median ratio),
%   aggregated/microbench.csv (per-config means), tables/throughput_parity.md.

\subsection{RQ3 --- Promotion rate and decision trail}
% Table: T1 columns plateau_events / proposals / promotions /
%   schema_valid_rate / fallback_rate

% =====================================================================
\section{Case Study: Plateau Breakthrough on cJSON}
\label{sec:case}
% TODO Day 6. Source: results/paper01_ai_recipe_mutation/tables/case_study/
%   (run picked by aggregate.py via 5-criterion filter, see
%   docs/papers/paper01_experiments_runbook.md §2.6).
%   Walk: blackboard at proposal time -> agent JSON -> recipe ops ->
%   micro-campaign reward -> main-run coverage delta in next 10 min.
%   Inline Fig 6 timeline (TikZ).
% Page budget: 0.5

% =====================================================================
\section{Related Work}
\label{sec:related}
% TODO Day 4. Build from T3 (related-work matrix). Source:
%   results/paper01_ai_recipe_mutation/tables/T3_related_work.md
% Lead with the "off the hot path" differentiator (only FuzzPilot row
% with "N" in the first column). Then validation, then audit log.
% Page budget: 0.5

% TODO Day 8: convert T3.md to LaTeX booktabs and \input{tables/T3_related_work.tex}
\begin{table}[t]
  \centering
  \footnotesize
  \fbox{\parbox{0.95\textwidth}{\centering\textit{[T3: Related-work matrix --- placeholder. Source: \texttt{tables/T3\_related\_work.md}]}}}
  \caption{Comparison across LLM-augmented and LLM-adjacent fuzzers.
  FuzzPilot is the only row with the model off the mutation hot path
  and an empirical validation step before promotion.}
  \label{tab:related}
\end{table}

% =====================================================================
\section{Limitations, Threats to Validity, and Future Work}
\label{sec:limits}
% TODO Day 4. Source: paper01_arxiv_placeholder.md
%   §"Honest limitations to surface" + §"Threats to validity".
%   Direct paste then polish. Critical items:
%   - Single target family (cJSON); libpng smoke as single extra point
%   - N=5 repeats; no formal significance tests
%   - No crash dedup pipeline yet
%   - Deferred ablations (random-recipe, ai-direct) -> Paper 2
%   - macOS/arm64 dev only; canonical Linux/x86_64
%   - microbench "vanilla" is havoc-cost proxy, not real AFL++ havoc
%   - LLM nondeterminism: temperature 0.0 + report schema_valid /
%     fallback rates instead of model identity claims
% Page budget: 0.25

% =====================================================================
\section{Conclusion}
\label{sec:conclusion}
% TODO Day 7: one paragraph echoing abstract + signposting Paper 2/3.
% Page budget: 0.25

% =====================================================================
% Reproducibility footnote (cited from §6 and §1)
% TODO Day 9: fill in:
%   - GitHub commit hash
%   - Git tag paper01-arxiv-v1
%   - AFL++ commit (v4.21c)
%   - Docker image digest
%   - Artifact tarball SHA-256
%   - Model name + temperature (deepseek-chat, 0.0)

\bibliographystyle{plain}
% TODO Day 4-8: build refs.bib alongside §8 prose. Citation TODOs in
%   tables/T3_related_work.md.
\bibliography{refs}

\end{document}
